%% ============================================================================
%% Chapter 2 — Literature Survey
%% AI Test Agent: Autonomous AI-Based Web Application Testing
%%                and Intelligent Bug Reporting
%% ============================================================================

\chapter{Literature Survey}
\label{chap:litsurvey}

This chapter presents a comprehensive review of the literature that underpins the
design and development of the AI Test Agent platform. The survey spans classical
software testing theory, modern automated testing frameworks, the emergence of
artificial intelligence and large language models (LLMs) in software engineering,
agentic AI paradigms, browser automation technologies, intelligent bug detection
research, and microservices-based software architecture. A comparative analysis
of existing tools and approaches is provided, culminating in an identification of
the research gap that the AI Test Agent addresses.

%%=============================================================================
\section{Introduction to Software Testing Research}
\label{sec:testing_research_intro}
%%=============================================================================

Software testing is a systematic process of evaluating a software system to
detect differences between existing and required conditions and to assess the
features of the software~\cite{myers2011art}. It is universally acknowledged as
one of the most critical activities in the software development life cycle
(SDLC), consuming an estimated 30--50\% of total project effort in
industrial practice~\cite{aniche2022testing}. The International Software
Testing Qualifications Board (ISTQB) defines testing as the process consisting
of all life cycle activities, both static and dynamic, concerned with planning,
preparation, and evaluation of a component or system and related work products
to determine that they satisfy specified requirements~\cite{istqb2023}.

Historically, software testing research has evolved through several distinct
phases. Early work in the 1970s and 1980s established the theoretical
foundations of test adequacy criteria, including statement coverage, branch
coverage, and path coverage~\cite{myers2011art}. The 1990s and 2000s saw the
maturation of automated test execution frameworks and the formalization of
model-based testing. The 2010s brought data-driven approaches, with machine
learning being applied to test case prioritization, fault prediction, and
automated test generation. Most recently, the 2020s have witnessed a
paradigm shift with the introduction of large language models capable of
understanding and generating both natural language and source
code~\cite{openai2023gpt4, chen2021codex}.

The AI Test Agent project is situated at this modern intersection of software
testing, agentic AI systems, and LLM-driven reasoning. It draws upon decades
of testing theory while leveraging the latest advances in AI to create an
autonomous testing platform that observes, reasons about, plans, and executes
functional tests on web applications without manual test script authoring.

%%=============================================================================
\section{Traditional Software Testing Techniques}
\label{sec:traditional_testing}
%%=============================================================================

Traditional software testing techniques form the bedrock upon which all
subsequent automation and AI-driven approaches are built. These techniques
are broadly categorized by the degree of internal knowledge available to
the tester.

\subsection{Black-Box Testing}
\label{subsec:black_box}

Black-box testing, also called specification-based testing, treats the system
under test as an opaque entity and derives test cases entirely from the
external specification~\cite{myers2011art}. The tester has no knowledge of or
access to the internal code structure. Key black-box techniques include:

\begin{itemize}[nosep]
  \item \textbf{Equivalence Partitioning:} The input domain is divided into
        classes of data from which test cases are derived, under the
        assumption that all members of a partition are processed equivalently.
  \item \textbf{Boundary Value Analysis:} Test cases are designed at the
        edges of equivalence partitions, as experience shows that errors
        tend to cluster at boundaries.
  \item \textbf{Decision Table Testing:} Combinatorial logic involving
        multiple conditions and their corresponding actions is captured
        in a decision table, from which test cases are systematically derived.
  \item \textbf{State Transition Testing:} The system is modeled as a finite
        state machine, and tests exercise transitions between states triggered
        by events and guarded by conditions.
\end{itemize}

The AI Test Agent employs a form of black-box testing in its core workflow:
the exploration service crawls the target application's user interface to
discover navigable elements, forms, buttons, and links. The LLM-based test
planner then generates functional test cases based solely on the observed
UI structure and inferred application semantics, without access to the
target application's source code.

\subsection{White-Box Testing}
\label{subsec:white_box}

White-box testing, or structural testing, leverages full knowledge of the
internal implementation to design test cases that exercise specific code
paths~\cite{myers2011art}. Techniques include statement coverage, branch
coverage, condition coverage, and path coverage. While white-box testing
achieves thorough code-level verification, it requires access to source
code and deep understanding of the implementation, making it impractical
for external or third-party application testing scenarios.

\subsection{Gray-Box Testing}
\label{subsec:gray_box}

Gray-box testing combines elements of both black-box and white-box
approaches. The tester possesses partial knowledge of the internal
structure---such as data flow, architectural patterns, or API
contracts---while designing tests that are executed from the external
interface. The AI Test Agent can be considered a gray-box approach in
the sense that the exploration service extracts structural information
(DOM elements, form structures, navigation patterns) from the
target application, providing the AI planner with partial internal
knowledge beyond what a pure black-box tester would possess.

\subsection{Regression Testing}
\label{subsec:regression}

Regression testing ensures that previously working functionality continues
to operate correctly after code changes, bug fixes, or feature
additions~\cite{istqb2023}. It is one of the most resource-intensive
testing activities, as the regression suite tends to grow monotonically
over the life of a project. Test case prioritization and selection
techniques have been extensively studied to manage regression testing
costs~\cite{aniche2022testing}. The AI Test Agent contributes to this
domain by enabling rapid regeneration of test suites for target
applications---since tests are generated autonomously from the current
UI state, the system naturally adapts to interface changes without
manual test maintenance.

%%=============================================================================
\section{Automated Functional Testing}
\label{sec:automated_testing}
%%=============================================================================

Automated functional testing involves the use of specialized software
tools to control the execution of tests and the comparison of actual
outcomes with predicted outcomes. This section reviews the major
frameworks and paradigms in web application test automation.

\subsection{Selenium}
\label{subsec:selenium}

Selenium is the most widely adopted open-source web testing
framework~\cite{selenium2004}. It provides the WebDriver API, which
enables programmatic control of web browsers (Chrome, Firefox, Edge,
Safari) through language bindings in Java, Python, C\#, JavaScript,
and Ruby. Selenium's architecture is based on a client--server model
where the WebDriver client sends commands to a browser-specific driver
binary, which translates them into browser-native automation protocols.

While Selenium revolutionized web testing, it has well-documented
limitations: flaky tests due to implicit timing dependencies, lack of
built-in auto-waiting mechanisms, difficulty handling modern single-page
application (SPA) frameworks, limited support for intercepting network
requests, and the overhead of managing separate driver binaries for each
browser version.

\subsection{Cypress}
\label{subsec:cypress}

Cypress emerged as a developer-centric alternative to Selenium,
executing tests directly within the browser's JavaScript runtime rather
than over a remote protocol. This in-process architecture provides
deterministic execution, automatic waiting, time-travel debugging, and
real-time reloading. However, Cypress is limited to Chromium-based
browsers and Firefox, does not support multi-tab or multi-domain testing
in a single test, and is restricted to the JavaScript/TypeScript
ecosystem.

\subsection{Playwright}
\label{subsec:playwright_overview}

Microsoft Playwright represents the latest generation of browser
automation libraries~\cite{playwright2024}. It supports Chromium,
Firefox, and WebKit from a single API, provides built-in auto-waiting,
network interception, and multiple browser context isolation. Playwright
is available in JavaScript/TypeScript, Python, Java, and .NET. Its
architecture communicates with browsers over the Chrome DevTools Protocol
(CDP) and equivalent protocols for Firefox and WebKit, enabling
fine-grained control over browser behavior. The AI Test Agent uses
Playwright for Java as its deterministic browser actuator, as discussed
further in Section~\ref{sec:playwright_detail}.

\subsection{Record-and-Playback vs.\ Script-Based Approaches}
\label{subsec:record_playback}

Record-and-playback tools (e.g., Selenium IDE, Playwright Codegen)
capture user interactions and replay them as automated tests. While
accessible to non-programmers, these tools produce brittle, hard-to-maintain
scripts that break with minor UI changes. Script-based approaches, where
engineers manually author test code using framework APIs, offer greater
flexibility and maintainability but require significant programming
expertise and time investment. The AI Test Agent transcends both paradigms:
it autonomously generates test plans from observed UI structure using LLM
reasoning, then executes them through Playwright's programmatic API---combining
the accessibility of record-and-playback with the robustness of
script-based execution.

%%=============================================================================
\section{Artificial Intelligence in Software Testing}
\label{sec:ai_testing}
%%=============================================================================

The application of artificial intelligence techniques to software testing
has been an active research area for over two decades, with significant
acceleration in recent years.

\subsection{Machine Learning for Test Case Generation}
\label{subsec:ml_test_gen}

Machine learning (ML) models have been applied to generate test cases
from various software artifacts. Supervised learning approaches train
models on historical test data to predict effective test inputs.
Reinforcement learning agents have been used to explore application
state spaces, learning optimal sequences of UI interactions that
maximize code coverage or fault detection~\cite{aniche2022testing}.
Search-based software testing (SBST) employs evolutionary algorithms
(genetic algorithms, particle swarm optimization) to evolve test suites
that optimize coverage criteria. Tools such as EvoSuite for Java unit
testing exemplify this approach, generating JUnit test suites that
achieve high branch coverage through guided search.

\subsection{Test Case Prioritization}
\label{subsec:test_prioritization}

Test case prioritization reorders test suites to maximize the rate of
fault detection, enabling earlier discovery of defects in time-constrained
testing environments~\cite{istqb2023}. ML-based prioritization
techniques use features such as code change history, test execution
history, code complexity metrics, and requirement traceability to predict
which tests are most likely to reveal faults. Deep learning models,
including recurrent neural networks, have been applied to learn temporal
patterns in test failure histories for improved prioritization accuracy.

\subsection{Fault Localization}
\label{subsec:fault_localization}

Fault localization identifies the most probable locations of defects in
source code, given a set of failing and passing test cases.
Spectrum-based fault localization (SBFL) techniques, such as Tarantula,
Ochiai, and DStar, rank program statements by suspiciousness scores
computed from test execution spectra. ML-enhanced fault localization
models incorporate additional features---such as code complexity,
developer change patterns, and code churn---to improve localization
accuracy. The AI Test Agent's bug analysis service applies a
complementary approach: rather than analyzing source code spectra, it
uses LLM-based reasoning over execution traces, failure screenshots, and
assertion results to generate root-cause hypotheses in natural language.

\subsection{Mutation Testing}
\label{subsec:mutation_testing}

Mutation testing evaluates test suite quality by introducing small
syntactic changes (mutants) into the source code and measuring the
proportion of mutants detected (killed) by the test
suite~\cite{aniche2022testing}. While computationally expensive, mutation
testing provides a rigorous measure of test effectiveness. AI techniques
have been applied to reduce mutation testing costs through equivalent
mutant detection and selective mutation strategies.

%%=============================================================================
\section{Large Language Models in Software Engineering}
\label{sec:llm_se}
%%=============================================================================

The emergence of large language models (LLMs) has fundamentally
transformed multiple aspects of software engineering, from code
generation to testing to documentation.

\subsection{Transformer Architecture and the GPT Family}
\label{subsec:transformers}

The transformer architecture, introduced by Vaswani et
al.~\cite{vaswani2017attention}, replaced recurrent neural networks
with self-attention mechanisms that process input sequences in parallel.
This architecture enabled the training of models on vastly larger
datasets, leading to emergent capabilities in language understanding
and generation. The Generative Pre-trained Transformer (GPT) family
of models, culminating in GPT-4~\cite{openai2023gpt4}, demonstrated
state-of-the-art performance across a wide range of natural language
processing and code-related tasks. GPT-4 is a multimodal model capable
of accepting both text and image inputs, producing text outputs with
human-level performance on numerous professional and academic
benchmarks.

\subsection{Code Generation with LLMs}
\label{subsec:code_gen}

OpenAI Codex, a GPT model fine-tuned on publicly available code
repositories, demonstrated that LLMs can generate functionally correct
code from natural language descriptions~\cite{chen2021codex}. Chen et
al.\ evaluated Codex on the HumanEval benchmark, a set of 164 Python
programming problems, achieving 28.8\% pass@1 accuracy with the initial
model, rising to 72.3\% with 100 samples (pass@100). Subsequent models
(GPT-3.5-turbo, GPT-4) substantially improved these numbers. The
capability to generate code from natural language is foundational to
the AI Test Agent's approach: the LLM receives structured descriptions
of discovered web pages and generates executable test plans expressed
as sequences of browser actions.

\subsection{Test Generation with LLMs}
\label{subsec:llm_test_gen}

Recent research has explored the use of LLMs for automated test
generation~\cite{king2024llmtest}. Studies demonstrate that LLMs can
generate unit tests, integration tests, and end-to-end tests from
source code, documentation, or UI descriptions. Key findings include:

\begin{itemize}[nosep]
  \item LLMs generate tests with reasonable code coverage, though they
        may produce tests with incorrect assertions (hallucinated oracles).
  \item Prompt engineering---including few-shot examples, chain-of-thought
        reasoning, and structured output formatting---significantly
        improves test quality.
  \item LLM-generated tests are most effective when combined with
        deterministic execution and validation, as the AI Test Agent
        does by separating LLM-based planning from Playwright-based
        execution.
\end{itemize}

The AI Test Agent builds on this research direction by using LLMs not
merely to generate isolated test functions but to reason about
end-to-end user workflows across multiple pages, producing structured
JSON test plans that are validated and executed deterministically.

%%=============================================================================
\section{Agentic AI Systems}
\label{sec:agentic_ai}
%%=============================================================================

Agentic AI refers to systems where an AI model operates as an
autonomous agent that perceives its environment, reasons about
observations, plans actions, and executes those actions in a feedback
loop~\cite{yang2023agentic, wang2024survey}. This paradigm extends
beyond single-query LLM usage into multi-step, goal-directed behavior.

\subsection{Autonomous Agents and Cognitive Architectures}
\label{subsec:autonomous_agents}

Wang et al.~\cite{wang2024survey} present a comprehensive survey on
LLM-based autonomous agents, identifying four key modules: profiling,
memory, planning, and action. The profiling module defines the agent's
role and expertise. The memory module maintains short-term (context
window) and long-term (external storage) information. The planning
module decomposes complex tasks into sub-goals through strategies such
as chain-of-thought reasoning, task decomposition, and self-reflection.
The action module translates plans into concrete operations in the
external environment.

\subsection{Observe--Reason--Act Loops}
\label{subsec:ora_loop}

The observe--reason--act (ORA) loop is the fundamental execution cycle
of an agentic AI system. At each iteration, the agent:

\begin{enumerate}[nosep]
  \item \textbf{Observes:} Perceives the current state of the
        environment through sensors, APIs, or tool outputs.
  \item \textbf{Reasons:} Applies internal models (LLM inference,
        retrieval-augmented generation, chain-of-thought) to interpret
        observations and decide on the next course of action.
  \item \textbf{Acts:} Executes the selected action using available
        tools, APIs, or actuators, thereby modifying the environment.
\end{enumerate}

The AI Test Agent implements a structured variant of this loop:
\texttt{OBSERVE} $\rightarrow$ \texttt{REASON} $\rightarrow$
\texttt{PLAN} $\rightarrow$ \texttt{ACT} $\rightarrow$
\texttt{OBSERVE~RESULT} $\rightarrow$ \texttt{ANALYZE}. The
exploration service observes the target application's UI structure.
The test/AI service reasons about what workflows and edge cases
merit testing and plans structured test cases. The execution service
acts by running these test cases through Playwright. The bug analysis
service observes the execution results and analyzes failures to
produce actionable bug reports.

\subsection{Tool-Using Agents}
\label{subsec:tool_agents}

A defining characteristic of modern agentic AI systems is their
ability to use external tools~\cite{wang2024survey}. Rather than
relying solely on the knowledge embedded in the LLM's parameters,
tool-using agents invoke APIs, execute code, query databases, browse
the web, and control external applications. In the AI Test Agent, the
LLM functions as the ``brain'' that generates plans and analyses,
while Playwright serves as the deterministic ``actuator'' that
executes browser operations. This strict separation---where the LLM
never directly executes arbitrary browser code---ensures
reproducibility, safety, and auditability of the testing process.

%%=============================================================================
\section{Browser Automation and Playwright}
\label{sec:playwright_detail}
%%=============================================================================

Browser automation is the technological foundation that enables
automated functional testing of web applications. This section provides
a detailed comparison of major browser automation technologies and
justifies the selection of Playwright for the AI Test Agent.

\subsection{Evolution of Browser Automation}
\label{subsec:browser_evolution}

Browser automation has evolved through several generations:

\begin{enumerate}[nosep]
  \item \textbf{First Generation (2004--2011):} Selenium RC, which
        injected JavaScript into the browser to simulate user actions.
        Limited by same-origin policy restrictions and timing
        instabilities.
  \item \textbf{Second Generation (2011--2018):} Selenium WebDriver,
        which introduced a standardized browser driver protocol (later
        formalized as the W3C WebDriver specification). The
        client--driver--browser architecture became the industry
        standard~\cite{selenium2004}.
  \item \textbf{Third Generation (2018--present):} CDP-based tools
        (Puppeteer, Playwright) that communicate directly with browser
        internals over the Chrome DevTools Protocol, enabling
        capabilities such as network interception, geolocation emulation,
        and device emulation~\cite{playwright2024}.
\end{enumerate}

\subsection{Comparison of Selenium, Cypress, and Playwright}
\label{subsec:framework_comparison}

Table~\ref{tab:browser_comparison} provides a feature-level comparison
of the three major browser automation frameworks.

\begin{table}[H]
\centering
\caption{Comparison of Browser Automation Frameworks}
\label{tab:browser_comparison}
\small
\begin{tabularx}{\textwidth}{@{} l X X X @{}}
\toprule
\textbf{Feature} & \textbf{Selenium} & \textbf{Cypress} & \textbf{Playwright} \\
\midrule
Multi-browser support & Chrome, Firefox, Edge, Safari & Chromium, Firefox (limited) & Chromium, Firefox, WebKit \\
Language bindings & Java, Python, C\#, JS, Ruby & JavaScript/TypeScript only & Java, Python, .NET, JS/TS \\
Auto-waiting & Manual (explicit/implicit waits) & Built-in & Built-in \\
Network interception & Via proxy (BrowserMob) & Built-in & Built-in \\
Multi-tab/multi-context & Limited & Not supported & Full support \\
Parallel execution & Via Selenium Grid & Limited & Native support \\
Mobile emulation & Via Appium (separate) & Limited & Built-in device emulation \\
iFrame handling & Supported (verbose) & Supported (limited) & Seamless \\
\bottomrule
\end{tabularx}
\end{table}

\subsection{Advantages of Playwright for AI Test Agent}
\label{subsec:playwright_advantages}

Playwright was selected as the browser automation engine for the AI
Test Agent for several technical reasons~\cite{playwright2024}:

\begin{itemize}[nosep]
  \item \textbf{Java Integration:} Playwright for Java integrates
        seamlessly with the Spring Boot microservices architecture,
        allowing browser automation to be embedded directly within the
        exploration and execution services.
  \item \textbf{Deterministic Auto-Waiting:} Playwright automatically
        waits for elements to be actionable (visible, stable, enabled)
        before performing operations, drastically reducing test
        flakiness compared to Selenium.
  \item \textbf{Browser Context Isolation:} Each test execution runs
        in an isolated browser context with independent cookies, storage,
        and cache, enabling parallel, interference-free test execution.
  \item \textbf{Network Interception:} Built-in route interception
        allows the system to monitor, mock, or modify network requests
        during test execution, which is valuable for testing error
        handling scenarios.
  \item \textbf{Screenshot and Trace Capture:} Native support for
        full-page screenshots and execution traces provides rich
        diagnostic data for the LLM-based bug analysis service.
\end{itemize}

%%=============================================================================
\section{Intelligent Bug Detection and Reporting}
\label{sec:intelligent_bug_detection}
%%=============================================================================

The automation of bug detection and reporting has been a long-standing
goal of software engineering research. Traditional approaches rely on
static analysis tools, dynamic analysis instruments, and manual
developer triage. Recent advances leverage AI to automate not only the
detection of failures but also the generation of structured,
human-readable bug reports.

\subsection{Automated Bug Reporting Research}
\label{subsec:automated_bug_reports}

Automated bug reporting systems aim to produce defect reports that
contain sufficient information for developers to understand, reproduce,
and fix the reported issue. Research has identified several essential
components of effective bug reports: a descriptive summary, steps to
reproduce, expected versus actual behavior, environmental context
(browser, OS, version), severity classification, and supporting
evidence (screenshots, logs, stack traces)~\cite{aniche2022testing}.

Traditional automated bug reporters extract information from crash
logs, exception traces, and execution recordings. However, they
typically lack the ability to interpret failure context semantically---that
is, to explain \emph{why} a failure occurred in terms of application
logic rather than merely recording \emph{what} happened at the
technical level.

\subsection{Root Cause Analysis with AI}
\label{subsec:root_cause_ai}

AI-driven root cause analysis (RCA) applies machine learning and,
more recently, LLM reasoning to identify the underlying cause of
observed failures. LLMs are particularly well-suited to this task
because they can process heterogeneous inputs---assertion failure
messages, HTML snapshots, screenshots, execution logs---and synthesize
a coherent natural-language explanation of the probable root
cause~\cite{openai2023gpt4}.

The AI Test Agent's bug analysis service implements LLM-based RCA
through the following pipeline:

\begin{enumerate}[nosep]
  \item Failed test execution results (including step-by-step outcomes,
        assertion messages, and screenshots) are collected by the
        execution service.
  \item These artifacts are forwarded to the bug analysis service, which
        invokes a LangGraph-orchestrated LLM workflow.
  \item The LLM analyzes the failure evidence, generates a structured
        bug report with severity classification (Critical, Major, Minor,
        Trivial) and a natural-language root-cause hypothesis, and
        returns the report for dashboard display.
\end{enumerate}

%%=============================================================================
\section{Microservices-Based Software Architecture}
\label{sec:microservices}
%%=============================================================================

The AI Test Agent is architected as a microservices-based system, a
design choice informed by extensive literature on modern distributed
software architecture.

\subsection{Microservices: Definition and Principles}
\label{subsec:microservices_def}

Fowler and Lewis define microservices as an architectural style that
structures an application as a collection of small, autonomous services,
each running in its own process, communicating over lightweight
mechanisms (typically HTTP REST APIs), and independently
deployable~\cite{fowler2014microservices}. Newman elaborates on the
organizational and technical principles: services are built around
business capabilities, own their data, can be deployed independently,
and embrace decentralized governance~\cite{newman2021microservices}.

\subsection{Key Microservices Patterns}
\label{subsec:microservices_patterns}

The AI Test Agent employs several established microservices patterns:

\begin{itemize}[nosep]
  \item \textbf{API Gateway Pattern:} A single entry point (Spring Cloud
        Gateway on port 8080) routes client requests to the appropriate
        backend service, handles cross-cutting concerns such as CORS
        configuration and JWT validation, and decouples the frontend
        from individual service endpoints~\cite{newman2021microservices}.
  \item \textbf{Service Registry and Discovery:} Netflix Eureka Server
        provides dynamic service registration and discovery. Each
        microservice registers itself with Eureka on startup, and the
        API Gateway uses load-balanced URIs (\texttt{lb://SERVICE-NAME})
        to route requests without hardcoded hostnames or
        ports~\cite{fowler2014microservices}.
  \item \textbf{Database-per-Service:} Each microservice owns a dedicated
        PostgreSQL database (\texttt{auth\_db}, \texttt{exploration\_db},
        \texttt{test\_db}, \texttt{execution\_db}, \texttt{bug\_db}).
        No service directly accesses another service's database; all
        inter-service communication occurs via REST
        APIs~\cite{newman2021microservices}.
  \item \textbf{Externalized Configuration:} Service configuration
        (database URLs, API keys, port assignments) is managed through
        environment variables and Spring Boot application properties,
        enabling deployment flexibility across development, testing,
        and production environments.
\end{itemize}

\subsection{Benefits of Microservices for the AI Test Agent}
\label{subsec:microservices_benefits}

The microservices architecture provides several specific advantages for
the AI Test Agent:

\begin{enumerate}[nosep]
  \item \textbf{Independent Scalability:} The execution service, which
        manages resource-intensive browser instances, can be scaled
        independently of the lightweight auth service.
  \item \textbf{Technology Heterogeneity:} While the backend services
        use Java/Spring Boot, the AI/agent layer can leverage Python-based
        frameworks (LangGraph, LangChain) where the ML ecosystem is
        richest.
  \item \textbf{Fault Isolation:} A failure in one service (e.g., the
        bug analysis service encountering an LLM API error) does not
        cascade to other services; the exploration and execution services
        continue to operate normally.
  \item \textbf{Team Autonomy:} The two-developer team was able to work
        on different services concurrently using feature branches, with
        well-defined REST API contracts as integration boundaries.
\end{enumerate}

%%=============================================================================
\section{Comparative Analysis of Existing Tools and Approaches}
\label{sec:comparative_analysis}
%%=============================================================================

Table~\ref{tab:comparative_analysis} presents a comparative analysis of
existing testing tools and approaches alongside the AI Test Agent,
evaluating them across key dimensions including autonomy, AI integration,
browser support, test generation methodology, bug reporting capability,
and architectural paradigm.

\begin{table}[H]
\centering
\caption{Comparative Analysis of Testing Approaches and Tools}
\label{tab:comparative_analysis}
\small
\renewcommand{\arraystretch}{1.25}
\begin{tabularx}{\textwidth}{@{} >{\raggedright\arraybackslash}p{2.1cm}
                                  >{\raggedright\arraybackslash}X
                                  >{\raggedright\arraybackslash}X
                                  >{\raggedright\arraybackslash}X
                                  >{\raggedright\arraybackslash}X
                                  >{\raggedright\arraybackslash}X
                                  >{\raggedright\arraybackslash}X @{}}
\toprule
\textbf{Criterion} &
\textbf{Selenium WebDriver} &
\textbf{Cypress} &
\textbf{Playwright} &
\textbf{SBST / EvoSuite} &
\textbf{LLM-Based Test Gen.} &
\textbf{AI Test Agent} \\
\midrule

Test Generation &
Manual script writing &
Manual script writing &
Manual or codegen &
Automated (search-based) &
Automated (LLM prompt) &
Fully autonomous (agentic LLM) \\

Target Scope &
Web UI &
Web UI &
Web UI &
Java unit tests &
Unit/API tests &
Web UI (end-to-end) \\

AI Integration &
None &
None &
None &
Evolutionary algorithms &
LLM for generation &
LLM for planning, analysis \& bug reporting \\

Browser Support &
Chrome, Firefox, Edge, Safari &
Chromium, Firefox &
Chromium, Firefox, WebKit &
N/A &
Varies &
Chromium, Firefox, WebKit (via Playwright) \\

Bug Reporting &
Manual by tester &
Manual by tester &
Trace viewer (manual) &
Not applicable &
Not integrated &
Automated LLM-generated reports \\

Architecture &
Monolithic client &
Monolithic runner &
Library/CLI &
CLI tool &
Script/notebook &
Microservices (6 services + gateway) \\

Autonomy Level &
None (fully manual) &
None (fully manual) &
Low (codegen assist) &
Medium (automated gen.) &
Medium (requires prompts) &
High (observe--reason--act loop) \\

Maintenance &
High (brittle selectors) &
Medium &
Low (auto-wait) &
Low (re-generate) &
Low (re-prompt) &
Minimal (re-explore \& regenerate) \\

\bottomrule
\end{tabularx}
\end{table}

As illustrated in Table~\ref{tab:comparative_analysis}, existing tools
occupy specific points in the automation spectrum. Selenium, Cypress,
and Playwright provide powerful browser automation capabilities but
require manual test authoring. Search-based approaches like EvoSuite
automate test generation but are limited to unit-level testing of Java
applications. LLM-based test generation research demonstrates promising
results but typically operates in isolated, single-shot prompting
scenarios without integrated execution or bug reporting. The AI Test
Agent uniquely combines autonomous exploration, agentic LLM-driven
test planning, deterministic Playwright execution, and intelligent
bug analysis within a cohesive microservices architecture.

%%=============================================================================
\section{Research Gap}
\label{sec:research_gap}
%%=============================================================================

The literature review reveals several significant gaps in existing
software testing tools and research that the AI Test Agent addresses:

\begin{enumerate}[nosep]
  \item \textbf{Fragmented Workflows:} Existing approaches treat test
        generation, test execution, and bug reporting as disconnected
        activities. A tester must separately write tests (or prompt an
        LLM), configure and run a test framework, manually analyze
        failures, and compose bug reports. The AI Test Agent integrates
        all four activities into a single, cohesive, end-to-end pipeline.

  \item \textbf{Lack of Autonomous Exploration:} Traditional test
        automation frameworks require testers to know, a priori, the
        structure of the application under test and the specific
        scenarios to verify. No mainstream tool autonomously explores
        a target application's UI to discover testable workflows. The
        AI Test Agent's exploration service fills this gap by crawling
        the target application and building structured page models that
        feed directly into AI-driven test planning.

  \item \textbf{No Closed-Loop Agentic Testing:} While agentic AI
        research has demonstrated observe--reason--act loops for
        general-purpose tasks~\cite{wang2024survey, yang2023agentic},
        the application of this paradigm specifically to functional
        web application testing---with deterministic execution
        guarantees and structured bug output---is not addressed by
        existing tools or research prototypes.

  \item \textbf{LLM Hallucination in Test Oracles:} LLM-generated
        tests are prone to hallucinated assertions---test oracles that
        do not correspond to actual application behavior. The AI Test
        Agent mitigates this by separating LLM planning (which produces
        structured action sequences) from deterministic Playwright
        execution (which performs actual assertions against the live
        application state), ensuring that test verdicts are grounded
        in real browser observations.

  \item \textbf{Manual Bug Triage:} Even in highly automated testing
        environments, failure analysis and bug report authoring remain
        manual, time-consuming activities. Existing tools generate raw
        logs and screenshots but do not interpret failures or suggest
        root causes. The AI Test Agent's LLM-powered bug analysis
        service generates structured, severity-classified bug reports
        with root-cause hypotheses, significantly reducing the manual
        effort required for defect triage.

  \item \textbf{Monolithic Tool Architecture:} Most testing tools are
        monolithic applications or CLI-based utilities that cannot be
        independently scaled, extended, or deployed. The AI Test Agent's
        microservices architecture enables independent scaling of
        resource-intensive components (browser execution, LLM inference),
        fault isolation, and technology heterogeneity---addressing
        the operational limitations of monolithic testing
        tools~\cite{fowler2014microservices, newman2021microservices}.
\end{enumerate}

\noindent
In summary, the AI Test Agent addresses a confluence of gaps at the
intersection of browser automation, agentic AI, LLM-driven reasoning,
and microservices architecture. By unifying autonomous exploration,
intelligent test planning, deterministic execution, and AI-powered bug
analysis into a single platform, it advances the state of practice
beyond what any individual existing tool or research approach provides.
